% preamble/colors.tex
% ============================================
% 配色方案：青蓝 & 金色 学术风格
% 主色：深青蓝（PhysTeal）—— 定理、命题、章节彩带
% 辅色：暖金（Gold）—— 装饰线、点缀
% ============================================

% ---- 青蓝色系（主结构色）----
\definecolor{PhysTeal}{RGB}{14, 120, 130}
\definecolor{PhysCyan}{RGB}{18, 155, 168}
\definecolor{PhysDark}{RGB}{6, 68, 74}

\colorlet{WaveMain}{PhysTeal}        % 主青蓝色 —— 封面侧边条、章节彩带
\colorlet{WaveDark}{PhysDark}        % 深青蓝色 —— 标题、强调
\colorlet{WaveLight}{PhysCyan}       % 浅青蓝色 —— 次强调

% ---- 暖金色系（装饰点缀用，与其它笔记一致）----
\definecolor{GoldDark}{RGB}{140, 100, 35}      % 暗金
\definecolor{GoldMain}{RGB}{178, 132, 55}      % 主金 —— 章节扉页装饰条、分割线
\definecolor{GoldLight}{RGB}{210, 172, 100}    % 浅金
\definecolor{GoldPale}{RGB}{242, 230, 205}     % 淡金背景

% ---- 兼容旧名称（各 preamble 文件内部桥接）----
\colorlet{BlueMain}{WaveMain}
\colorlet{BlueDark}{WaveDark}
\colorlet{BlueLight}{WaveLight}
\colorlet{GreenMain}{WaveMain}
\colorlet{OrangeMain}{WaveMain}       % chapter.title.tex 中的彩带

% ---- 旧几何色名桥接（消除 geometry.tex 残留依赖）----
\colorlet{mainblue}{WaveMain}
\colorlet{mediumblue}{WaveDark}
\colorlet{lightblue}{WaveLight}
\colorlet{grayblue}{WaveMain!50}
\colorlet{brightblue}{GoldMain}
\colorlet{paleblue}{WaveMain!5}

% ---- 代码块样式 ----
\lstset{
	backgroundcolor=\color{WaveMain!4},
	basicstyle=\ttfamily\small,
	keywordstyle=\color{WaveMain},
	commentstyle=\color{WaveMain!40},
	escapeinside={(*@}{@*)},
}
